\PassOptionsToPackage{dvipsnames}{xcolor}
\documentclass[ignorenonframetext, hyperref=unicode,compress]{beamer}


\usepackage{cmap}
\usepackage[T2A]{fontenc}
\usepackage[utf8x]{inputenc}
\usepackage[bulgarian]{babel}
\selectlanguage{bulgarian}

\usepackage{graphicx}
\usepackage{listings}

\hypersetup{
	colorlinks=true,
	linkcolor=blue,
	filecolor=blue,
	urlcolor=blue,
	anchorcolor=blue,
	citecolor=blue
}

\lstset{language=java, 
  numbers=left, 
  numberstyle=\tiny,
  stepnumber=1, 
  numbersep=3pt, 
  tabsize=2, 
  texcl,
  basicstyle=\ttfamily\small,
  identifierstyle=\ttfamily\small,
  keywordstyle=\sffamily\bfseries\small,
  extendedchars=true, inputencoding=utf8x,
  backgroundcolor=\color[rgb]{1,1,0.845},
  escapeinside={/*@}{@*/}}

\newcommand{\Cpp}{{\ttfamily\bfseries C++}}
\newcommand{\CC}{{\ttfamily\bfseries C}}

% \usepackage{algpseudocode}

%\usepackage{ucs}

%\logo{\includegraphics[height=0.4cm]{../macros/logo_elsys.png}}

\titlegraphic{\href{http://creativecommons.org/licenses/by-sa/3.0/}{\includegraphics{pics/cc.png}}}

\newcommand{\authors}{%
\author[Л.~Чорбаджиев]{Любомир Чорбаджиев}
\institute[ELSYS] {%
Технологично училище ``Електронни системи'' \\
Технически университет, София
}}

\usepackage{tikz}
\usetikzlibrary{arrows,shapes.misc,calc,backgrounds,fit,decorations.pathmorphing,positioning,chains,scopes,topaths,automata,shapes,trees}

\tikzset{obj/.style={rectangle,minimum size=6mm,
				very thick,draw=red!50!black!50,top color=white,bottom color=red!50!black!20, font=\ttfamily},
ref/.style={rectangle,
		minimum size=8mm,
		% The rest
		very thick,draw=black!50,
		top color=white,bottom color=black!20},
	every node/.style={font=\ttfamily},
	>=stealth',
	thick,
	draw=black!50,
	shorten >=1pt,node distance=2cm and 2.5cm,on grid}


\mode<article>
{

}

\mode<presentation>
{
  \usetheme[secheader=true]{Madrid}
  \usecolortheme{crane}
  \usefonttheme[onlylarge]{structurebold}
  \setbeamercovered{transparent}
}

\usepackage[unicode]{hyperref}

%%% Local Variables: 
%%% mode: latex
%%% TeX-master: t
%%% End: 


\title{Обекти в Java}
\authors
\date{\today}


\begin{document}

\frame{\titlepage}

\begin{frame}
\small
{\bf Забележка:} Тази лекция е адаптация на:
\begin{itemize}
  \item \href{http://ocw.mit.edu/NR/rdonlyres/Electrical-Engineering-and-Computer-Science/6-092January--IAP--2006/97145060-6EC3-44D9-9B9B-35B412CCAE45/0/lecture1b.pdf}{Corey McCaffrey: {\em Java Objects}} from
\href{http://ocw.mit.edu/OcwWeb/Electrical-Engineering-and-Computer-Science/6-092January--IAP--2006/CourseHome/index.htm}{
{\em 6.092: Java for 6.170} (MIT OpenCourseWare:
Massachusetts Institute of Technology)}\\
{\bf Лиценз:}
\href{http://ocw.mit.edu/OcwWeb/web/terms/terms/index.htm\#cc}{Creative commons
BY-NC-SA}  

\end{itemize}

\end{frame}

\begin{frame}
\frametitle{Съдържание}
\tableofcontents %[hideallsubsections]
\end{frame}



\begin{frame}[containsverbatim]\frametitle{Препратки}
\begin{itemize}
 \item Когато в Java дефинирате променлива от даден клас, то всъщност дефинирате препратка (reference). 
Дефинирането на променлива не създава обект.
 \item Препратката или сочи към обект или е \lstinline{null}.
\end{itemize}

\begin{lstlisting}
Integer i1=null;
Integer i2=new Integer(3);
\end{lstlisting}
\end{frame}

\begin{frame}[containsverbatim]\frametitle{Обекти}
\begin{itemize}
 \item Обектът е инстанция (екземпляр) на даден клас.
 \item За да се създаде обект е необходимо да се използва оператора \lstinline{new}. 
Този оператор извиква конструктор на обекта.
\end{itemize}

\begin{lstlisting}
Integer i1=null;
Integer i2=new Integer(3);
\end{lstlisting}
\end{frame}


\begin{frame}[containsverbatim]\frametitle{Пример}
\begin{lstlisting}
public  class AssignmentReview {
 	public static void main(String[] args) {
		Integer num=null;
 		num = new Integer(3);
		Integer x = num;
 		Integer y = new Integer(3);
		Integer z;
	}
}
\end{lstlisting}
\begin{center}
 
\begin{tikzpicture}

\node (int1) [obj] at (0,0) {Integer:3};
\node (num) [ref] at (-0.5,2) {num};
\node (x)   [ref] at (0.5,2) {x};

\path[->] (num) edge node{} (int1)
       		(x) edge (int1);

\node (int2) [obj] at (3,0) {Integer:3};
\node (y) [ref] at (3,2) {y};

\path[->] (y) edge node{} (int2);

\node (null) [obj] at (6,0) {null};
\node (z) [ref] at (6,2) {z};

\path[->] (z) edge node{} (null);
\end{tikzpicture}
\end{center}

\end{frame}

\begin{frame}[containsverbatim]\frametitle{Препратки \lstinline{null}}
\begin{itemize}
 \item Ненасочените препратки сочат към \lstinline{null}
 \item \lstinline{null} не е обект (няма полета, нито методи)
\end{itemize}

\begin{lstlisting}
Integer z;
z.intValue(); //този ред ще генерира NullPointerException
\end{lstlisting}
\end{frame}

\begin{frame}[containsverbatim]\frametitle{Обекти и препратки}
\begin{itemize}
 \item Операторът за присвояване (\lstinline{=}) насочва препратката към обект
 \item Методите са в състояние да променят вътрешното състояние на обекта
 \item Може да имате няколко препратки към един и същ обект, така че се пазете от странични ефекти
\end{itemize}
\end{frame}

\begin{frame}[containsverbatim]\frametitle{Обекти и препратки}
\begin{lstlisting}
public  class MutationExample {
 	public static void main(String[] args) {
		List<String> a = new ArrayList<String>();
		List<String> b = a; 
		a.add("Hello, world!");
		System.out.println(b);
	}
}
\end{lstlisting}

\begin{center}
\begin{tikzpicture}

\node (list) [obj]  at (0,0) {List<String>[''Hello, World!'']};
\node (a)  [ref] at (-1,2) {a};
\node (b)  [ref] at (1,2) {b};

\path[->] (a) edge node{} (list)
       		(b) edge (list);

\end{tikzpicture}
\end{center}

\end{frame}

\begin{frame}[containsverbatim]\frametitle{Статични и нестатични членове}
\begin{itemize}
 \item Методите и полетата могат да се декларират като статични
 \item Статичните методи/полета принадлежат на класа
 \item Нестатичните методи/полета принадлежат на обекта
\end{itemize}
\end{frame}

\begin{frame}[containsverbatim]\frametitle{Пример: Нестатично поле}
\begin{lstlisting}
public  class Bean {
 	public int beanCounter = 0;

 	public Bean() {
 		beanCounter++;
	}

 	public static void main(String[] args) {
 		new Bean();
 		new Bean();
 		Bean bean = new Bean();
 		System.out.println(bean.beanCounter);
			//prints 1
	}
}
\end{lstlisting}
\end{frame}


\begin{frame}[containsverbatim]\frametitle{Пример: Статично поле}
\begin{lstlisting}
public  class Bean {
 	public static int beanCounter = 0;

 	public Bean() {
 		beanCounter++;
	}

 	public static void main(String[] args) {
 		new Bean();
 		new Bean();
 		Bean bean = new Bean();
 		System.out.println(bean.beanCounter);
			//prints 3
	}
}
\end{lstlisting}
\end{frame}

\begin{frame}[containsverbatim]\frametitle{Пример: Нестатичен метод}
\begin{lstlisting}
public  class Bean {
 	private boolean planted = false;

 	public void plantBean() {
 		bean.planted = true;
	}

 	public static void main(String[] args) {
 		Bean bean = new Bean();
 		bean.plantBean(); // Invoked on instance
	}
}
\end{lstlisting}
\end{frame}

\begin{frame}[containsverbatim]\frametitle{Пример: Статичен метод}
\begin{lstlisting}
public  class Bean {
 	private boolean planted = false;

 	public static void plantBean(Bean bean) {
 		bean.planted = true;
	}

 	public static void main(String[] args) {
 		Bean bean = new Bean();
 		Bean.plantBean(bean); // Invoked on class
 		bean.plantBean(bean); // Legal but inadvisable!
	}
}
\end{lstlisting}
\end{frame}

\begin{frame}[containsverbatim]\frametitle{Пример: Обекти, предавани с препратки}
\begin{lstlisting}
public  static <T> void removeFirst(List<T> list) {
		list.remove(0);
}

public  static void main(String[] args) {
 	List<String> myList= new ArrayList<String>();
	myList.add("Cat"); 
	myList.add("Dog");
 	System.out.println(myList); // Prints [Cat,Dog]
	removeFirst(myList);
 	System.out.println(myList); // Prints [Dog]
}
\end{lstlisting}
\end{frame}

\begin{frame}[containsverbatim]\frametitle{Обекти и препратки}

\begin{center}
\begin{tikzpicture}

\node (list) [obj]  at (0,0) {List<String>[''Cat'',''Dog'']};
\node (myList)  [ref] at (-1,2) {myList};
\node (lst)  [ref] at (1,2) {list};

\path[->] (myList) edge node{} (list)
       		(lst) edge (list);

\end{tikzpicture}
\end{center}

\end{frame}

\begin{frame}[containsverbatim]\frametitle{Блокове и видимост}
\begin{itemize}
 \item Блокът (\{ \}) определя областта на видимост на препратките
 \item Препратките съществуват от момента на създаването им до момента, 
в който излязат от областта на видимост
 \item Полетата могат да се използват в рамките на класа
 \item Параметрите на даден метод могат да се използват в рамките на метода
\end{itemize}
\end{frame}

\begin{frame}[containsverbatim]\frametitle{Пример: Област на видимост}
\begin{lstlisting}
public  class ScopeExample {
 	private int field;

 	public int method(int parameter) {
 		int localVar1;
 		if (field > 0) {
 			int x;
		}
 		int localVar2;
	}
}
\end{lstlisting}
\end{frame}

\begin{frame}[containsverbatim]\frametitle{Пример: Област на видимост}
\begin{lstlisting}
public  class ScopeExample {
 	private int field;

 	public int method(int parameter) {
 		int field; // Legal, but hides field!
 		int localVar;
 		if (this.field > 0) { // Accesses field
 			int x;
		}
 		int localVar; // Illegal: same scope
	}
}
\end{lstlisting}
\end{frame}

\begin{frame}[containsverbatim]\frametitle{Изводи}
\begin{itemize}
 \item Препратките сочат към обекти. Трябва да се прави разлика между препратка и обект.
 \item Пазете се от препратки, които са \lstinline{null}
 \item Не извиквайте статични методи чрез инстанции
 \item Ако не искате даден обект да се променя, то когато го предавате като аргумент трябва изрично да направите копие.
 \item Минимизирайте видимостта на променливите
\end{itemize}
\end{frame}



\end{document}
